\documentclass[11pt]{article}

\usepackage[T1]{fontenc}
\usepackage[margin=1in]{geometry}
\usepackage{amsmath,amssymb,amsthm}
\usepackage{hyperref}
\usepackage{enumitem}
\usepackage{booktabs}

\hypersetup{
  colorlinks=true,
  linkcolor=blue,
  urlcolor=blue
}

\newcommand{\fl}{\operatorname{fl}}
\newcommand{\Frob}[1]{\lVert #1\rVert_F}
\newcommand{\Prob}{\mathbb{P}}
\newcommand{\E}{\mathbb{E}}
\newcommand{\R}{\mathbb{R}}
\newcommand{\Nat}{\mathbb{N}}
\newcommand{\lean}[1]{\texttt{\detokenize{#1}}}

\newtheorem{theorem}{Theorem}
\newtheorem{corollary}[theorem]{Corollary}
\newtheorem{lemma}[theorem]{Lemma}

\theoremstyle{definition}
\newtheorem{definition}[theorem]{Definition}
\newtheorem{assumption}[theorem]{Assumptions}

\theoremstyle{remark}
\newtheorem*{leanref}{Lean reference}
\newtheorem*{remark}{Remark}

\title{Theorem and Corollary Summary of the Lean Proof\\
for RandNLA Algorithm 1 Stability}
\author{LeanFpAnalysis repository}
\date{}

\begin{document}
\sloppy
\maketitle

\section*{Overview}

This note summarizes the Lean formalization of the entrywise floating-point
stability analysis for Algorithm 1, the element-wise sampling meta-algorithm
from the RandNLA paper. The formalized sampling probabilities are the
squared-magnitude probabilities
\[
  p_{ij} =
  \frac{A_{ij}^2}{\sum_{k,l} A_{kl}^2}
  =
  \frac{A_{ij}^2}{\Frob{A}^2}.
\]

The document is written as mathematical theorem and corollary statements. Each
statement also records the corresponding Lean theorem or definition name.

\section*{Notation and assumptions}

Let \(A \in \R^{m \times n}\), let \(s\) be the sampling parameter from
Algorithm 1, and fix an entry \((i,j)\). Let
\(\widetilde A^{(0)}\) be the initial sketch and
\(\widehat{\widetilde A}\) be the floating-point sketch produced by the
algorithm.

The floating-point model is the repository's standard relative-error model:
for each elementary operation,
\[
  \fl(x \circ y) = (x \circ y)(1+\delta),
  \qquad |\delta| \le u,
\]
where \(u\) is the unit roundoff. The accumulated-error notation is Higham's
\[
  \gamma_q = \frac{q u}{1-q u}.
\]
Lean represents the required side condition \(q u < 1\) as
\(\lean{gammaValid fp q}\).

Throughout the main entrywise theorem, the relevant assumptions are:
\[
  s \ne 0,
  \qquad
  A_{ij} \ne 0,
  \qquad
  \lean{gammaValid fp q_{ij}},
  \qquad
  \lean{gammaValid fp (q_{ij}+1)}.
\]

\section*{Deterministic arithmetic}

\begin{definition}[Squared-magnitude increment]
For a sampled nonzero entry \((i,j)\), Algorithm 1 adds
\[
  \frac{A_{ij}}{s p_{ij}}
\]
to the sketch entry. Under squared-magnitude probabilities this increment is
\[
  \frac{A_{ij}}{s p_{ij}}
  =
  \frac{\Frob{A}^2}{s A_{ij}}.
\]
\end{definition}

\begin{leanref}
Formalized as \(\lean{elementwiseIncrement_sqMag_eq}\).
\end{leanref}

\begin{theorem}[One-sample floating-point stability]
Fix a scalar sketch value \(x\). Suppose \(s \ne 0\) and \(A_{ij} \ne 0\).
The floating-point update
\[
  \fl\left(x + \fl\left(\frac{A_{ij}}{s p_{ij}}\right)\right)
\]
satisfies
\[
\begin{aligned}
&\left|
  \widehat x -
  \left(x + \frac{\Frob{A}^2}{s A_{ij}}\right)
\right| \\
&\qquad\le
  \left|x + \frac{\Frob{A}^2}{s A_{ij}}\right|u
  +
  \left|\frac{\Frob{A}^2}{s A_{ij}}\right|u(1+u).
\end{aligned}
\]
\end{theorem}

\begin{leanref}
Formalized as \(\lean{fl_elementwiseUpdateEntry_sqMag_error_bound}\).
\end{leanref}

\section*{Trace reduction}

\begin{definition}[Hit counter]
For a deterministic trace
\[
  \texttt{samples : Fin steps -> Fin m x Fin n},
\]
the hit counter \(q_{ij}\) is the number of trace steps that sample entry
\((i,j)\):
\[
  q_{ij} =
  \#\{t : \texttt{Fin steps} \mid \texttt{samples}(t)=(i,j)\}.
\]
\end{definition}

\begin{leanref}
Formalized as \(\lean{hitCount samples i j}\).
\end{leanref}

\begin{lemma}[Trace equals repeated scalar updates]
For a fixed output entry \((i,j)\), the exact trace and the floating-point trace
depend on the full sample trace only through \(q_{ij}\). In other words, the
entrywise trace update is equivalent to repeating the corresponding scalar
entry update exactly \(q_{ij}\) times.
\end{lemma}

\begin{leanref}
Formalized as
\[
\begin{array}{l}
  \lean{elementwiseTraceSketch_eq_repeat_of_hitCount},\\
  \lean{fl_elementwiseTraceSketch_eq_repeat_of_hitCount}.
\end{array}
\]
\end{leanref}

\section*{Main deterministic result}

\begin{theorem}[Deterministic entrywise stability for Algorithm 1]
Fix an entry \((i,j)\), and let \(q_{ij}\) be its hit count in the trace. Under
the assumptions above, Algorithm 1 with squared-magnitude sampling satisfies
\[
\begin{aligned}
&\left|
  \widehat{\widetilde A}_{ij}
  -
  \left(
    \widetilde A^{(0)}_{ij}
    +
    q_{ij}\frac{\Frob{A}^2}{s A_{ij}}
  \right)
\right| \\
&\qquad\le
  \left|\widetilde A^{(0)}_{ij}\right|\gamma_{q_{ij}}
  +
  q_{ij}
  \left|\frac{\Frob{A}^2}{s A_{ij}}\right|
  \gamma_{q_{ij}+1}.
\end{aligned}
\]
\end{theorem}

\begin{leanref}
Formalized as \(\lean{fl_elementwiseTraceSketch_sqMag_error_bound}\).
\end{leanref}

\begin{corollary}[Comparison with the exact trace]
The same deterministic bound can be written as a direct comparison between the
floating-point trace and the exact arithmetic trace produced by the same sample
sequence.
\end{corollary}

\begin{leanref}
Formalized as \(\lean{fl_elementwiseTraceSketch_sqMag_error_bound_exact}\).
\end{leanref}

\begin{corollary}[Zero-initial-sketch Algorithm 1]
If Algorithm 1 starts from the zero sketch, so that
\(\widetilde A^{(0)}_{ij}=0\), then
\[
\left|
  \widehat{\widetilde A}_{ij}
  -
  q_{ij}\frac{\Frob{A}^2}{s A_{ij}}
\right|
\le
q_{ij}
\left|\frac{\Frob{A}^2}{s A_{ij}}\right|
\gamma_{q_{ij}+1}.
\]
\end{corollary}

\begin{leanref}
Formalized as \(\lean{fl_elementwiseSketch_zero_init_sqMag_error_bound}\).
\end{leanref}

\section*{Deterministic budget from a count bound}

\begin{definition}[Stability budget]
For a deterministic upper bound \(Q\) on the hit counter, define
\[
\begin{aligned}
B(Q)
&=
  \left|\widetilde A^{(0)}_{ij}\right|\gamma_Q
  +
  Q
  \left|\frac{\Frob{A}^2}{s A_{ij}}\right|
  \gamma_{Q+1}.
\end{aligned}
\]
\end{definition}

\begin{leanref}
Formalized as \(\lean{sqMagTraceErrorBudget}\).
\end{leanref}

\begin{lemma}[Budget monotonicity]
If \(q_{ij} \le Q\), then the deterministic stability budget at
\(q_{ij}\) is bounded by the deterministic stability budget at \(Q\):
\[
  B(q_{ij}) \le B(Q).
\]
\end{lemma}

\begin{leanref}
Formalized as \(\lean{sqMagTraceErrorBudget_mono}\).
\end{leanref}

\begin{corollary}[Stability from a deterministic count bound]
If \(q_{ij} \le Q\), then Algorithm 1 satisfies the deterministic entrywise
stability bound with \(B(Q)\) in place of the trace-specific budget
\(B(q_{ij})\).
\end{corollary}

\begin{leanref}
Formalized as
\[
  \lean{fl_elementwiseTraceSketch_sqMag_error_bound_of_hitCount_le}.
\]
\end{leanref}

\section*{Random trace model}

\begin{definition}[Random trace]
Let \(X\) be a random trace in a finite probability space. For fixed
\((i,j)\), the random variable of interest is
\[
  q_{ij}(X) = \lean{hitCount (X omega) i j}.
\]
For fixed \(A\), \(s\), floating-point model, and initial sketch, this hit
counter is the only random quantity in the deterministic stability budget.
\end{definition}

\begin{leanref}
The finite probability interface is \(\lean{FiniteProbability}\), and the
identity writing the hit count as a sum of indicators is
\(\lean{hitCount_eq_sum_indicator}\).
\end{leanref}

\begin{theorem}[Mean hit count under the marginal sampling law]
Assume that each sample step hits entry \((i,j)\) with marginal probability
\(p_{ij}\). Then
\[
  \E[q_{ij}] = \texttt{steps}\cdot p_{ij}.
\]
For squared-magnitude sampling,
\[
  p_{ij} = \frac{A_{ij}^2}{\Frob{A}^2}.
\]
\end{theorem}

\begin{leanref}
Formalized as
\[
  \lean{expectationNat_hitCount_eq_steps_mul_hitProb}
  \quad\text{and its squared-magnitude specialization.}
\]
\end{leanref}

\section*{Markov high-probability result}

\begin{theorem}[Markov concentration for the hit counter]
Assume only the marginal law
\[
  \Prob(\text{step } t \text{ hits } (i,j)) = p_{ij}
\]
for every step \(t\). Then, for any natural number \(Q\),
\[
  \Prob(q_{ij} \le Q)
  \ge
  1 - \frac{\texttt{steps}\cdot p_{ij}}{Q+1}.
\]
Equivalently, if
\[
  \frac{\texttt{steps}\cdot p_{ij}}{Q+1} \le \delta,
\]
then
\[
  \Prob(q_{ij} \le Q) \ge 1-\delta.
\]
\end{theorem}

\begin{leanref}
Formalized as
\[
  \lean{hitCount_concentration_sqMag_markov}
  \quad\text{and}\quad
  \lean{hitCount_concentrates_sqMag}.
\]
\end{leanref}

\begin{corollary}[Chosen Markov budget]
Let \(\delta>0\), and choose
\[
  Q_M =
  \left\lceil
    \frac{\texttt{steps}\cdot p_{ij}}{\delta}
  \right\rceil.
\]
Then
\[
  \Prob(q_{ij} \le Q_M) \ge 1-\delta.
\]
\end{corollary}

\begin{leanref}
The chosen count budget is formalized as
\(\lean{sqMagMarkovHitCountBudget}\). The tail inequality is formalized as
\(\lean{sqMagMarkovHitCountBudget_tail}\).
\end{leanref}

\begin{corollary}[High-probability stability with the Markov budget]
Under the squared-magnitude marginal sampling law, with \(Q_M\) as above,
Algorithm 1 satisfies the deterministic stability bound with budget \(B(Q_M)\)
with probability at least \(1-\delta\):
\[
  \Prob\left(
    \left|
      \widehat{\widetilde A}_{ij}
      -
      \left(
        \widetilde A^{(0)}_{ij}
        +
        q_{ij}\frac{\Frob{A}^2}{s A_{ij}}
      \right)
    \right|
    \le B(Q_M)
  \right)
  \ge 1-\delta.
\]
\end{corollary}

\begin{leanref}
Formalized as
\(\lean{highProbability_sqMagTraceStability_of_markov_budget}\).
\end{leanref}

\section*{Pairwise-Chebyshev high-probability result}

\begin{theorem}[Centered second moment under pairwise independence]
Assume the marginal law above and assume that distinct step-hit indicators for
entry \((i,j)\) are pairwise independent. Then Lean proves an explicit formula
for
\[
  M =
  \E\left[
    \left(q_{ij}-\texttt{steps}\cdot p_{ij}\right)^2
  \right].
\]
This is the variance-scale quantity used by Chebyshev's inequality.
\end{theorem}

\begin{leanref}
The centered moment is formalized as
\(\lean{hitCountPairwiseCenteredMoment}\). The theorem identifying it with the
expectation is
\[
  \lean{expectationReal_hitCount_centered_sq_eq_pairwise}.
\]
\end{leanref}

\begin{theorem}[Chebyshev concentration around the mean]
Under the same marginal and pairwise-independence assumptions, if
\[
  \frac{M}{\varepsilon^2} \le \delta,
\]
then
\[
  \Prob\left(
    \left|
      q_{ij} - \texttt{steps}\cdot p_{ij}
    \right|
    \le \varepsilon
  \right)
  \ge 1-\delta.
\]
For squared-magnitude sampling, \(p_{ij}=A_{ij}^2/\Frob{A}^2\).
\end{theorem}

\begin{leanref}
Formalized as
\(\lean{hitCount_concentrates_sqMag_around_mean_pairwise}\).
\end{leanref}

\begin{corollary}[Chosen Chebyshev budget]
Let \(\delta>0\), let
\[
  M =
  \E\left[
    \left(q_{ij}-\texttt{steps}\cdot p_{ij}\right)^2
  \right],
\]
and choose the conservative radius
\[
  \varepsilon_C = \frac{M}{\delta}+1.
\]
Then Lean proves
\[
  \frac{M}{\varepsilon_C^2} \le \delta.
\]
Consequently, if
\[
  Q_C =
  \left\lceil
    \texttt{steps}\cdot p_{ij}+\varepsilon_C
  \right\rceil,
\]
then
\[
  \Prob(q_{ij} \le Q_C) \ge 1-\delta.
\]
\end{corollary}

\begin{leanref}
The chosen radius and budget are formalized as
\[
  \lean{chebyshevHitCountRadius},
  \quad
  \lean{sqMagChebyshevHitCountBudget}.
\]
The key arithmetic lemmas are
\[
\begin{array}{l}
  \lean{hitCountPairwiseCenteredMoment_div_chebyshevRadius_sq_le},\\
  \lean{sqMagChebyshevHitCountBudget_mean_add_radius_le}.
\end{array}
\]
\end{leanref}

\begin{corollary}[High-probability stability with the Chebyshev budget]
Under the squared-magnitude marginal law and pairwise independence of distinct
step-hit indicators, with \(Q_C\) as above, Algorithm 1 satisfies the
deterministic stability bound with budget \(B(Q_C)\) with probability at least
\(1-\delta\):
\[
  \Prob\left(
    \left|
      \widehat{\widetilde A}_{ij}
      -
      \left(
        \widetilde A^{(0)}_{ij}
        +
        q_{ij}\frac{\Frob{A}^2}{s A_{ij}}
      \right)
    \right|
    \le B(Q_C)
  \right)
  \ge 1-\delta.
\]
\end{corollary}

\begin{leanref}
Formalized as
\[
  \lean{highProbability_sqMagTraceStability_of_pairwise_chebyshev_budget}.
\]
\end{leanref}

\section*{Chernoff high-probability result}

\begin{theorem}[Canonical product-sampler MGF]
Assume \(\Frob{A}^2>0\), and let \(P_A\) be the finite product probability
law that samples every trace step independently with probability
\[
  p_{kl}=\frac{A_{kl}^2}{\Frob{A}^2}.
\]
Then, for every \(\lambda\),
\[
  \E_{P_A}[\exp(\lambda q_{ij})]
  =
  \left(1+p_{ij}(\exp(\lambda)-1)\right)^{\texttt{steps}}.
\]
In particular, for \(\lambda>0\),
\[
  \E_{P_A}[\exp(\lambda q_{ij})]
  \le
  \exp\left(
    \texttt{steps}\cdot p_{ij}\cdot(\exp(\lambda)-1)
  \right).
\]
\end{theorem}

\begin{leanref}
The product law is \(\lean{sqMagTraceProbability}\). The exact MGF identity is
\[
  \lean{sqMagTraceProbability_expectationReal_exp_hitCount_eq},
\]
and the derived Chernoff MGF bound is
\[
  \lean{sqMagTraceProbability_chernoff_mgf_bound}.
\]
\end{leanref}

\begin{corollary}[Chernoff concentration for the canonical sampler]
Under the same canonical independent sampling law, for every natural number
\(Q\) and every \(\lambda>0\),
\[
  \Prob_{P_A}(q_{ij}\le Q)
  \ge
  1 -
  \exp\left(
    \texttt{steps}\cdot p_{ij}\cdot(\exp(\lambda)-1)
    -\lambda(Q+1)
  \right).
\]
\end{corollary}

\begin{leanref}
The generic MGF condition is still named
\(\lean{sqMagHitCountChernoffMGFBound}\), but the canonical-sampler theorem
uses the proved product-law MGF internally:
\[
  \lean{hitCount_concentrates_sqMag_chernoff_independent}.
\]
\end{leanref}

\begin{corollary}[Chosen fixed-parameter Chernoff budget]
Let \(\delta>0\) and \(\lambda>0\). For the canonical independent sampler,
choose
\[
  Q_{\mathrm{Ch}}(\lambda)
  =
  \left\lceil
    \frac{
      \texttt{steps}\cdot p_{ij}\cdot(\exp(\lambda)-1)-\log(\delta)
    }{\lambda}
  \right\rceil.
\]
Then the Chernoff tail is at most \(\delta\), and hence
\[
  \Prob_{P_A}(q_{ij}\le Q_{\mathrm{Ch}}(\lambda)) \ge 1-\delta.
\]
\end{corollary}

\begin{leanref}
The chosen budget is \(\lean{sqMagChernoffHitCountBudget}\), and the tail
inequality is \(\lean{sqMagChernoffHitCountBudget_tail}\). The fully derived
canonical-sampler concentration statement is
\[
  \lean{hitCount_concentrates_sqMag_chernoff_independent_budget}.
\]
\end{leanref}

\begin{corollary}[Canonical high-probability stability with the Chernoff budget]
For the canonical independent squared-magnitude sampler, Algorithm 1 satisfies
the deterministic stability bound with budget \(B(Q_{\mathrm{Ch}}(\lambda))\)
with probability at least \(1-\delta\). That is,
\[
  \Prob_{P_A}\left(
    \left|
      \widehat{\widetilde A}_{ij}
      -
      \left(
        \widetilde A^{(0)}_{ij}
        +
        q_{ij}\frac{\Frob{A}^2}{s A_{ij}}
      \right)
    \right|
    \le B(Q_{\mathrm{Ch}}(\lambda))
  \right)
  \ge 1-\delta.
\]
\end{corollary}

\begin{leanref}
Formalized as
\[
  \lean{highProbability_sqMagTraceStability_of_independent_chernoff_budget}.
\]
\end{leanref}

\begin{theorem}[Optimized Chernoff exponent for the canonical sampler]
For a supplied natural budget \(Q\), assume
\[
  0 < \mu = \texttt{steps}\cdot p_{ij}
  \quad\text{and}\quad
  \mu < Q+1.
\]
Lean chooses the optimizing parameter
\[
  \lambda_* = \log\left(\frac{Q+1}{\mu}\right)
\]
and proves the corresponding optimized Chernoff tail. Thus, if
\[
  \texttt{sqMagChernoffOptimizedHitCountTail}(Q) \le \delta,
\]
then
\[
  \Prob_{P_A}(q_{ij}\le Q) \ge 1-\delta.
\]
\end{theorem}

\begin{leanref}
Formalized as
\[
\begin{array}{l}
  \lean{hitCount_concentrates_sqMag_chernoff_optimized_independent},\\
  \lean{highProbability_sqMagTraceStability_of_independent_chernoff_optimized_tail_budget}.
\end{array}
\]
\end{leanref}

\section*{Generic probability-to-stability transfer}

\begin{theorem}[Transferring hit-count concentration to stability]
Let \(Q\) be any deterministic natural number. If
\[
  \Prob(q_{ij} \le Q) \ge 1-\delta,
\]
then the Algorithm 1 stability event with deterministic budget \(B(Q)\) also
has probability at least \(1-\delta\):
\[
  \Prob(\text{stability bound with budget }B(Q)) \ge 1-\delta.
\]
\end{theorem}

\begin{leanref}
Formalized as
\[
  \lean{highProbability_sqMagTraceStability_of_hitCount_concentration}.
\]
\end{leanref}

\begin{corollary}[Tail-budget version]
If the Markov tail condition
\[
  \frac{\texttt{steps}\cdot p_{ij}}{Q+1} \le \delta
\]
holds, then the stability event with budget \(B(Q)\) holds with probability at
least \(1-\delta\).
\end{corollary}

\begin{leanref}
Formalized as
\[
  \lean{highProbability_sqMagTraceStability_of_marginal_hitProb_of_tail_budget}.
\]
\end{leanref}

\begin{corollary}[Deviation-budget version]
If Chebyshev gives
\[
  \Prob\left(
    \left|q_{ij}-\texttt{steps}\cdot p_{ij}\right| \le \varepsilon
  \right)
  \ge 1-\delta
\]
and the chosen natural budget satisfies
\[
  \texttt{steps}\cdot p_{ij}+\varepsilon \le Q,
\]
then the stability event with budget \(B(Q)\) holds with probability at least
\(1-\delta\).
\end{corollary}

\begin{leanref}
Formalized as
\[
  \lean{highProbability_sqMagTraceStability_of_pairwise_hitCount_deviation}.
\]
\end{leanref}

\section*{What is, and is not, proved}

The formalization proves entrywise floating-point stability for Algorithm 1,
first deterministically and then with high probability by proving explicit
concentration bounds for the random hit counter \(q_{ij}\).

The Markov corollary needs only marginal hit probabilities, but it is crude.
The Chebyshev corollary uses pairwise independence and proves concentration
around the mean before transferring that result into the floating-point
stability bound. The Chernoff corollary constructs the canonical independent
squared-magnitude product law for Algorithm 1 traces, proves the hit-count MGF
bound from that law, and gives the exponentially sharper tail used in the
tightest counter-level analysis.

The current formalization does not yet prove matrix-level RandNLA accuracy
results such as spectral-norm concentration, unbiasedness of the full sketch,
or sample-complexity theorems.

\section*{Lean artifact table}

\begin{center}
\begin{tabular}{ll}
\toprule
Mathematical object or result & Lean artifact \\
\midrule
Rectangular Frobenius norm &
\lean{frobNormSqRect} \\
Squared-magnitude probabilities &
\lean{sqMagProb} \\
Squared-magnitude increment identity &
\lean{elementwiseIncrement_sqMag_eq} \\
One-sample stability &
\begin{tabular}[t]{@{}l@{}}
\lean{fl_elementwiseUpdateEntry}\\
\lean{_sqMag_error_bound}
\end{tabular} \\
Trace hit count &
\lean{hitCount} \\
Deterministic budget &
\lean{sqMagTraceErrorBudget} \\
Canonical product trace law &
\lean{sqMagTraceProbability} \\
Canonical Chernoff MGF theorem &
\begin{tabular}[t]{@{}l@{}}
\lean{sqMagTraceProbability}\\
\lean{_chernoff_mgf_bound}
\end{tabular} \\
Exact trace reduction &
\begin{tabular}[t]{@{}l@{}}
\lean{elementwiseTraceSketch}\\
\lean{_eq_repeat_of_hitCount}
\end{tabular} \\
Floating-point trace reduction &
\begin{tabular}[t]{@{}l@{}}
\lean{fl_elementwiseTraceSketch}\\
\lean{_eq_repeat_of_hitCount}
\end{tabular} \\
Main deterministic theorem &
\begin{tabular}[t]{@{}l@{}}
\lean{fl_elementwiseTraceSketch}\\
\lean{_sqMag_error_bound}
\end{tabular} \\
Zero-initial corollary &
\begin{tabular}[t]{@{}l@{}}
\lean{fl_elementwiseSketch}\\
\lean{_zero_init_sqMag_error_bound}
\end{tabular} \\
Markov chosen budget &
\lean{sqMagMarkovHitCountBudget} \\
Markov stability corollary &
\begin{tabular}[t]{@{}l@{}}
\lean{highProbability_sqMagTraceStability}\\
\lean{_of_markov_budget}
\end{tabular} \\
Chebyshev chosen budget &
\lean{sqMagChebyshevHitCountBudget} \\
Chebyshev stability corollary &
\begin{tabular}[t]{@{}l@{}}
\lean{highProbability_sqMagTraceStability}\\
\lean{_of_pairwise_chebyshev_budget}
\end{tabular} \\
Chernoff chosen budget &
\lean{sqMagChernoffHitCountBudget} \\
Optimized Chernoff tail &
\lean{sqMagChernoffOptimizedHitCountTail} \\
Chernoff stability corollary &
\begin{tabular}[t]{@{}l@{}}
\lean{highProbability_sqMagTraceStability}\\
\lean{_of_independent_chernoff_budget}
\end{tabular} \\
Optimized Chernoff stability &
\begin{tabular}[t]{@{}l@{}}
\lean{highProbability_sqMagTraceStability}\\
\lean{_of_independent_chernoff}\\
\lean{_optimized_tail_budget}
\end{tabular} \\
\bottomrule
\end{tabular}
\end{center}

The main source files are:
\[
\lean{LeanFpAnalysis/FP/Algorithms/RandNLA/ElementwiseSampling.lean}
\]
and
\[
\lean{LeanFpAnalysis/FP/Algorithms/RandNLA/HitCountConcentration.lean}.
\]

\end{document}
